\begin{proposition}[A two-prime norm obstruction for integral subcategories]
\label{prop:frontier-integral-center}
Let $\mathcal B$ be a fusion category with no nontrivial invertible simple
objects. Suppose that its integral-dimensional simple objects span a fusion
subcategory $\mathcal E$, and that the absolute field norm of
$\operatorname{FPdim}(\mathcal B)$ has at most two distinct prime divisors.
Then $\mathcal E=\mathrm{Vec}$.
\end{proposition}
\begin{proof}
Write $N=\operatorname{FPdim}(\mathcal E)\in\mathbb Z_{>0}$ and
$K=\mathbb Q(\operatorname{FPdim}(\mathcal B))$.
Subcategory divisibility \cite[Proposition~8.15]{ENO} implies that
$\operatorname{FPdim}(\mathcal B)/N$ is an algebraic integer.
Taking its norm shows that every prime divisor of $N$ divides the norm of
$\operatorname{FPdim}(\mathcal B)$. Thus $N=p^a q^b$ for at most two primes.
The category $\mathcal E$ is solvable by \cite[Theorem~1.6]{ENOSolvable}.
If nontrivial, it has a nontrivial invertible simple by
\cite[Proposition~4.5(iv)]{ENOSolvable}, contradicting the hypothesis on
$\mathcal B$.
\end{proof}

For any hypothetical category $\mathcal C$ with ring $R_{8,3,439}$ or
$R_{8,3,567}$, the already verified center-invertible obstruction gives
$\operatorname{Inv}Z(\mathcal C)=\{\mathbf1\}$.
Its simple dimensions are six copies of $1$ and two copies of
$d=3+2\sqrt3$, so $D=\operatorname{FPdim}(\mathcal C)=48+24\sqrt3$.
Every central simple restricts to an object of dimension $a+bd$ with
$a,b\in\mathbb Z_{\geq0}$. It has integral dimension precisely when $b=0$.
These objects form a fusion subcategory: their forgetful images lie in the
pointed source subcategory, which is closed under tensor products and duals.
Since $\operatorname{FPdim}Z(\mathcal C)=D^2$
\cite[Proposition~8.12]{ENO}, its integral subcategory dimension $N$ satisfies
\[
 \operatorname{Norm}_{\mathbb Q(\sqrt3)/\mathbb Q}(D^2/N)
 =\frac{576^2}{N^2}\in\mathbb Z,
 \qquad N\mid576=2^6 3^2.
\]
Proposition~\ref{prop:frontier-integral-center} therefore shows that the tensor
unit is the only integral-dimensional central simple. This conclusion requires
no pivotal or pseudo-unitary hypothesis.

The same center condition holds for the $R_{5,3,40}$ Morita dual forced by
the full-$S_3$ algebra in a hypothetical $R_{8,3,439}$ category. Its source
dimensions are $1,1,2,d,d$, and the integral source objects form the
$\operatorname{Rep}(S_3)$ based subcategory. The same proof removes every
nonunit central restriction column supported on $1,s,V$. This assertion is
about the specified Morita dual; it does not impose trivial center invertibles
on arbitrary categorifications of $R_{5,3,40}$.
